\documentclass[11pt,letterpaper]{article}

% ---------- Typography: serif body, sans headings (institutional memo standard) ----------
\usepackage[margin=1.0in,top=0.85in,bottom=0.85in]{geometry}
\usepackage[T1]{fontenc}
\usepackage{mathpazo}            % Palatino body, the classic finance-document serif
\usepackage{helvet}              % Helvetica for headings and labels only
\usepackage{microtype}           % refined kerning and justification
\usepackage{xcolor}
\usepackage[colorlinks=true,urlcolor=navy,linkcolor=navy]{hyperref}
\urlstyle{sf}                    % URLs in the document's sans font, not monospace
\usepackage{enumitem}
\usepackage{fancyhdr}
\usepackage{lastpage}
\usepackage{array}
\usepackage{graphicx}
\usepackage{booktabs}
\usepackage{needspace}           % keeps headings with the text that follows them
\usepackage{float}

% Matches the deck's macro name so the two documents' headline strings stay
% byte-for-byte comparable. No \bm here: the report's body text is not bold, so the
% plain sign is the one that matches its surroundings.
\newcommand{\apx}{\ensuremath{\approx}}

\definecolor{navy}{HTML}{1F3864}
\definecolor{midgray}{HTML}{595959}
\definecolor{hairline}{HTML}{C8C8C8}

\linespread{1.06}                % tightened from Week 5's 1.10 to hold the report to two pages
\setlength{\parindent}{0pt}
\setlength{\parskip}{3.2pt}

% Never strand a single line of a paragraph at the top or foot of a page.
\widowpenalty=10000
\clubpenalty=10000

% ---------- Footer: understated, informational ----------
\pagestyle{fancy}
\fancyhf{}
\renewcommand{\headrulewidth}{0pt}
\fancyfoot[L]{\sffamily\fontsize{7.5}{9}\selectfont\color{midgray}ENNACIRI \textbar\ PREDICTING U.S. BANK DISTRESS \textbar\ AUGUST 2026}
\fancyfoot[R]{\sffamily\fontsize{7.5}{9}\selectfont\color{midgray}\thepage\ OF \pageref{LastPage}}

% ---------- Section label: small uppercase bold sans, navy, corporate memo register ----------
\newcommand{\sectionhead}[1]{%
  \needspace{4\baselineskip}\vspace{4pt}%
  {\sffamily\bfseries\fontsize{9.5}{12}\selectfont\color{navy}\MakeUppercase{#1}}\par
  \vspace{-9pt}\textcolor{hairline}{\rule{\textwidth}{0.5pt}}\par\vspace{0pt}}

% Sub-label within a section
\newcommand{\sublabel}[1]{%
  \needspace{3\baselineskip}%
  \vspace{3pt}{\sffamily\bfseries\fontsize{9}{11}\selectfont\color{midgray}#1}\par\vspace{2pt}}

% Answer body
\newcommand{\answer}[1]{#1\par\vspace{3pt}}

% ---------- Storyboard machinery ----------
% The assignment is the headlines in order and nothing else, so the page has to let a
% reader run the headline-only test on it directly: these lines are the whole slide
% sequence, read top to bottom with nothing in between.
\newcommand{\phaselabel}[2]{%
  \needspace{4\baselineskip}\vspace{5pt}%
  {\sffamily\bfseries\fontsize{8}{10}\selectfont\color{navy}\MakeUppercase{#1}}%
  {\sffamily\fontsize{8}{10}\selectfont\color{midgray}\quad #2}\par\vspace{3pt}}

% #1 = slide number, #2 = the assertion headline
% The number sits in a fixed gutter and the headline in a minipage beside it, so a
% headline that wraps keeps its second line under the first rather than sliding back
% out to the margin.
\newlength{\slidegutter}\setlength{\slidegutter}{2.1em}
\newcommand{\slidehl}[2]{%
  \needspace{2\baselineskip}%
  \par\vspace{3pt}%
  \noindent\makebox[\slidegutter][l]{%
    \raisebox{0.02\baselineskip}{\sffamily\bfseries\fontsize{8.5}{11}\selectfont\color{hairline}#1}}%
  \begin{minipage}[t]{\dimexpr\textwidth-\slidegutter\relax}#2\end{minipage}\par}

\begin{document}

% ---------- Title block ----------
\vspace*{-14pt}
{\sffamily\fontsize{8.5}{11}\selectfont\color{midgray}MSDS 696 DATA SCIENCE PRACTICUM II \quad\textbar\quad WEEK 6 STATUS REPORT}\par\vspace{10pt}
{\sffamily\bfseries\fontsize{21}{25}\selectfont\color{navy}Predicting U.S. Bank Distress}\par\vspace{9pt}
{\sffamily\fontsize{9}{12}\selectfont Oussama Ennaciri \quad\textbar\quad Week 6 \quad\textbar\quad August 9, 2026}\par\vspace{5pt}
\textcolor{navy}{\rule{\textwidth}{1.2pt}}

% ---------- Project summary ----------
\sectionhead{Project summary}
\answer{This project builds an early-warning model that uses public FDIC and FRED quarterly data to flag which U.S.\ banks are at risk of becoming undercapitalized. The data consists of five linked FDIC tables (1.68M bank-quarter records of each bank's quarterly financial and regulatory filings) plus a FRED macro table. Modelling and the executive visuals are complete. This week storyboarded the final presentation: a one-sentence recommendation, then the slide headlines in order.}

% ---------- Milestones ----------
\sectionhead{Milestones}
\answer{\begin{enumerate}[topsep=1pt,itemsep=1pt,leftmargin=16pt,label=\textcolor{navy}{\bfseries\arabic*.}]
  \item Literature review of bank failure and existing early-warning models (\textbf{DONE})
  \item Collect the data (FDIC API, 1984 to present) (\textbf{DONE})
  \item Understand the data through exploratory data analysis (\textbf{DONE})
  \item Prepare the data for machine learning (\textbf{DONE})
  \item Train and compare four candidate models (\textbf{DONE})
  \item Move the critical-undercapitalization test to the Tier 1 capital basis and re-run the comparison (\textbf{DONE})
  \item Build the executive visuals (\textbf{DONE})
  \item Storyboard the final presentation (\textbf{DONE})
  \item Build the final deck and rehearse it (\textbf{IN PROGRESS})
\end{enumerate}}

% ---------- Last week's To Do ----------
\sectionhead{Last week's ``To Do''}
\answer{From the Week 5 report:
\begin{enumerate}[topsep=2pt,itemsep=1pt,leftmargin=16pt,label=\textcolor{navy}{\bfseries\arabic*.}]
  \item \textit{Tune gradient boosting on held-back training years, then use the test period once at the end.}
  \item \textit{Explore a second model, built on funding data, to cover the failures this one cannot see: the large banks that go through deposit runs rather than capital erosion.}
\end{enumerate}

\vspace{3pt}
\textbf{Both dropped.} The instructor asked me to stop adding models and spend the time left making the process clearer, which I am working on for the final presentation.}

% ---------- This week's progress ----------
% The heading, the intro line and the first headline have to travel together: left to
% itself the break stranded the section label at the foot of page 1 with its content
% overleaf. Seven baselines is heading + sub-label + intro line + the first headline.
\needspace{7\baselineskip}
\sectionhead{This week's progress}

\sublabel{The storyboard}
\answer{Ten slides, headlines in order.}

\slidehl{1}{Predicting U.S.\ Bank Distress}

\phaselabel{The answer}{}
\slidehl{2}{Supervisors should rank the review queue with a model instead of ranking on capital.}

\phaselabel{Situation}{}
\slidehl{3}{\apx4,900 banks file every quarter, so the review queue has to put the potentially troubled ones first.}
\slidehl{4}{The obvious way to rank that queue is by capital against assets, the measure the regulation itself uses.}

\phaselabel{Complication}{}
\slidehl{5}{Ranked by capital ratio, half the troubled banks never reach the top of the list.}

\phaselabel{Resolution}{}
\slidehl{6}{An alternative is to rank the queue with a machine learning model trained on past filings.}
\slidehl{7}{The model ranks the queue better than the capital ratio.}
\slidehl{8}{\apx60\% of the banks the model caught got a full year of warning.}
\slidehl{9}{Every troubled bank that went on to fail was flagged first.}
\slidehl{10}{Rank the review queue with the model, not the capital ratio.}

% ---------- Issues & discussion ----------
\sectionhead{Issues \& discussion}
\answer{\begin{itemize}[topsep=1pt,itemsep=4pt,leftmargin=14pt,label=\textcolor{navy}{\footnotesize\textbullet}]
  \item \textbf{The headline percentages do not stand on their own.} Each is measured under conditions the headline does not state, so a reader cannot tell what the numbers apply to. That context will come from the visualizations and the speaker notes.
  \item \textbf{The test set is thin.} 160 cases from 44 banks across nine years, against 161,117 filings.
\end{itemize}}

% ---------- Next week's To Do ----------
\sectionhead{Next week's ``To Do''}
\answer{Build the ten slides against these headlines and rehearse the five-minute version.}

% ---------- Resources ----------
\sectionhead{Resources}
\answer{{\fontsize{8.8}{10.8}\selectfont\setlength{\parskip}{2pt}\par\vspace{3pt}
\newcommand{\refentry}[1]{\hangindent=1.6em\hangafter=1 #1\par}
\refentry{Alley, M. (2013). \textit{The craft of scientific presentations: Critical steps to succeed and critical errors to avoid} (2nd ed.). Springer. (The assertion-evidence structure the storyboard follows.)}
\refentry{Minto, B. (2009). \textit{The pyramid principle: Logic in writing and thinking}. Pearson. (Situation--Complication--Resolution.)}
\refentry{Capital measures and capital category definitions, 12 C.F.R.\ \S\ 324.403 (the PCA thresholds behind the target label).}
}}

% ---------- LLM note ----------
\vspace{2pt}
{\small\itshape\color{midgray}\textbf{Note on LLM use:} This document was typeset in LaTeX using Claude Code, which handled the formatting, styling code, and wording polish. All content, research, and decisions are my own, and every suggested edit was reviewed and approved by me.}

\end{document}
